% --- Содержимое файла materials/4_5.tex ---

\subsection{Базисы в гильбертовых пространствах. Равенство Парсеваля}

\begin{definition}[Базис Шаудера]
	Пусть $E$ --- линейное нормированное пространство. Система $\{e_n\} \subset E$ называется \textbf{базисом Шаудера} в $E$, если для любого $x \in E$ существует \textit{единственный} набор чисел $\{\alpha_n\} \subset \mathbb{K}$ такой, что ряд сходится к $x$ по норме пространства:
	\[ x = \sum_{n = 1}^{\infty} \alpha_n e_n \]
\end{definition}

\begin{proposition}[Неравенство Бесселя]
	Пусть $E$ --- евклидово пространство, $\{e_i\}$ --- \textbf{ортонормированная система}. Тогда для любого $x \in E$:
	\[\sum_{k = 1}^{\infty}|(x, e_k)|^2 \leqslant \|x\|^2\]
\end{proposition}
\begin{proof}[Доказательство неравенства Бесселя]
	Достаточно заметить, что в силу ортонормированности системы $\{e_k\}$ выполнено равенство
	$$
	0 \le \left\|x-\sum_{k=1}^{n}(x,e_k)e_k\right\|^2
	= \|x\|^2-\sum_{k=1}^{n}|(x,e_k)|^2.
	$$
	
	Поскольку левая часть неотрицательна, то
	$$
	\sum_{k=1}^{n}|(x,e_k)|^2 \le \|x\|^2
	\quad \text{для любого } n\ge1.
	$$
	
	Переходя к пределу при $n\to\infty$, получаем
	$$
	\sum_{k=1}^{\infty}|(x,e_k)|^2 \le \|x\|^2.
	$$
\end{proof}

\begin{theorem}[Неравенство Коши--Буняковского]
	Пусть $E$ --- евклидово пространство. Тогда для любых $x,y\in E$ выполнено неравенство
	$$
	|(x,y)|\le \|x\|\,\|y\|.
	$$
\end{theorem}

\begin{proof}
	Случай $y=0$ тривиален, поэтому можно считать, что $y\ne 0$. Тогда система
	$$
	\left\{\frac{y}{\|y\|}\right\}
	$$
	является ортонормированной. По неравенству Бесселя для любого $x\in E$ имеем
	$$
	\left|\left(x,\frac{y}{\|y\|}\right)\right|^2 \le \|x\|^2.
	$$
	Но
	$$
	\left(x,\frac{y}{\|y\|}\right)=\frac{(x,y)}{\|y\|},
	$$
	следовательно,
	$$
	\frac{|(x,y)|^2}{\|y\|^2}\le \|x\|^2.
	$$
	Умножая на $\|y\|^2$ и извлекая квадратный корень, получаем
	$$
	|(x,y)|\le \|x\|\,\|y\|.
	$$
\end{proof}

\begin{proposition}[Экстремальное свойство коэффициентов Фурье]\label{minproperty}
	Пусть $\{e_k\}_{k=1}^n$ --- ортонормированная система в евклидовом пространстве $E$. Тогда для любого $x \in E$ и любых чисел $\alpha_k$:
	\[\left\|x - \sum_{k = 1}^n \alpha_ke_k\right\| \geqslant \left\|x - \sum_{k = 1}^n (x, e_k)e_k\right\|\]
	Равенство достигается тогда и только тогда, когда $\alpha_k = (x, e_k)$ (коэффициенты Фурье).
\end{proposition}

\begin{proof}
	[Доказательство для случая, когда $\mathbb{K}=\mathbb{R}$]
	В силу ортонормированности системы $\{e_k\}_{k=1}^n$ имеем
	$$
	\left\|x-\sum_{k=1}^n \alpha_k e_k\right\|^2
	=
	\left(
	x-\sum_{k=1}^n \alpha_k e_k,\;
	x-\sum_{k=1}^n \alpha_k e_k
	\right).
	$$
	
	Раскрывая скалярное произведение, получаем
	$$
	=
	\|x\|^2
	-2\sum_{k=1}^n \alpha_k (x,e_k)
	+\sum_{k=1}^n \alpha_k^2.
	$$
	
	С другой стороны,
	$$
	\left\|x-\sum_{k=1}^n (x,e_k)e_k\right\|^2
	=
	\|x\|^2-\sum_{k=1}^n (x,e_k)^2.
	$$
	
	Поэтому
	$$
	\left\|x-\sum_{k=1}^n \alpha_k e_k\right\|^2
	=
	\left\|x-\sum_{k=1}^n (x,e_k)e_k\right\|^2
	+\sum_{k=1}^n\big((x,e_k)-\alpha_k\big)^2.
	$$
	
	Отсюда следует неравенство
	$$
	\left\|x-\sum_{k=1}^n \alpha_k e_k\right\|
	\ge
	\left\|x-\sum_{k=1}^n (x,e_k)e_k\right\|.
	$$
	
	Причём равенство достигается тогда и только тогда, когда
	$$
	\alpha_k=(x,e_k),\quad k=1,\dots,n.
	$$
\end{proof}

\begin{idea}
	Это алгебраическое выражение геометрического факта: перпендикуляр короче наклонной. Проекция вектора $x$ на подпространство, натянутое на $e_k$, дается именно коэффициентами Фурье. Любой другой выбор коэффициентов $\alpha_k$ отклоняет нас от проекции.
\end{idea}

\begin{definition}[Сепарабельное пространство]\label{separability}
	Метрическое пространство $X$ называется \textbf{сепарабельным}, если в нем существует не более чем счетное всюду плотное множество.
	\[ \exists S \subset X : |S| \le \aleph_0 \quad \text{и} \quad \overline{S} = X \]
\end{definition}

\begin{idea}
	\textbf{Суть:} Сепарабельность означает, что пространство хоть и «большое» (несчетное), но его скелет — «маленький» (счетный). 
	Это позволяет нам контролировать пространство, используя последовательности (например, приближать любой элемент пределом последовательности из счетного набора).
	
	\textbf{Примеры:}
	\begin{itemize}
		\item $\mathbb{R}$ — сепарабельно (множество $\mathbb{Q}$ счетно и плотно).
		\item $C[a, b]$ — сепарабельно (многочлены с рациональными коэффициентами).
		\item $l_\infty$ — \textbf{не} сепарабельно (там слишком много «изолированных» точек).
	\end{itemize}
\end{idea}

\begin{theorem}[Критерий базиса]\label{thm4.4}
	Пусть $H$ --- сепарабельное гильбертово пространство, $\dim{H} = \infty$, и $e = \{e_n\}$ --- ортонормированная система. Следующие условия эквивалентны:
	\begin{enumerate}
		\item $e$ --- ортонормированный базис;
		\item Система $e$ полна (замкнутая линейная оболочка совпадает с $H$: $\overline{[e]} = H$);
		\item Для любого $h \in H$ выполнено \textbf{равенство Парсеваля}: $\|h\|^2 = \sum_{n = 1}^\infty |(h, e_n)|^2$;
		\item Аннулятор системы тривиален: $e^\perp = \{0\}$ (нет ненулевого вектора, ортогонального всем $e_n$).
	\end{enumerate}
\end{theorem}

\begin{proof}
	\begin{itemize}
		\item $(1) \Rightarrow (2)$
		По определению базиса, любой $x$ есть предел частичных сумм ряда $\sum \alpha_k e_k$, а частичные суммы лежат в $[e]$. Значит $x \in \overline{[e]}$.
		
		\item $(2) \Rightarrow (1)$
		\begin{enumerate}
			\item \textbf{Существование:} Зафиксируем $h \in H, \varepsilon > 0$. Так как $\overline{[e]} = H$, существует линейная комбинация $\sum_{k=1}^n \alpha_k e_k$, приближающая $h$ с точностью $\varepsilon$.
			\item По предложению \ref{minproperty}, сумма с коэффициентами Фурье приближает еще лучше:
			\[ \left\|h - \sum_{k=1}^n (h, e_k)e_k\right\| \leqslant \left\|h - \sum_{k=1}^n \alpha_k e_k\right\| < \varepsilon. \]
			\item Значит, ряд Фурье сходится к $h$.
			\item \textbf{Единственность:} Если $h = \sum \beta_n e_n$, то скалярно умножив на $e_k$ (и пользуясь непрерывностью скалярного произведения), получим $(h, e_k) = \sum \beta_n (e_n, e_k) = \beta_k$. Коэффициенты определены однозначно.
		\end{enumerate}
		
		\item $(1) \Leftrightarrow (3)$
		Ранее показано: $\left\|h - \sum_{k=1}^n (h, e_k)e_k\right\|^2 = \|h\|^2 - \sum_{k=1}^n |(h, e_k)|^2$.
		Сходимость ряда к $h$ (левая часть $\to 0$) равносильна сходимости суммы квадратов коэффициентов к квадрату нормы (правая часть $\to 0$).
		
		\item $(2) \Leftrightarrow (4)$
		По теореме о проекции, $H = \overline{[e]} \oplus \overline{[e]}^\perp$.
		Так как $e^\perp = [e]^\perp = \overline{[e]}^\perp$, то условие $\overline{[e]} = H$ равносильно $e^\perp = \{0\}$.
	\end{itemize}
\end{proof}

\begin{corollary}
	В любом бесконечномерном сепарабельном гильбертовом пространстве $H$ существует счетный ортонормированный базис.
\end{corollary}

\begin{idea}[Смысл теоремы]
	
	Эта теорема утверждает, что «хороший» базис в бесконечномерном пространстве ведет себя так же, как обычный базис $(\vec{i}, \vec{j}, \vec{k})$ в школьной геометрии.
	
	Все 4 условия говорят об одном и том же — \textbf{о полноте набора векторов}:
	
	\begin{itemize}
		
		\item \textbf{(1) Базис:} Любой вектор можно собрать из кирпичиков $e_n$.
		
		\item \textbf{(2) Полнота системы:} Линейными комбинациями $e_n$ можно подобраться сколь угодно близко к любой точке (нет недостижимых областей).
		
		\item \textbf{(3) Равенство Парсеваля:} Это \textbf{теорема Пифагора} для бесконечности. Квадрат длины вектора равен сумме квадратов его проекций на оси ($c^2 = a^2 + b^2 + \dots$).
		
		\item \textbf{(4) Тривиальный аннулятор:} Если вектор перпендикулярен всем осям координат сразу, то это может быть только нулевой вектор. Если бы нашелся ненулевой перпендикуляр, значит, мы «пропустили» какое-то измерение.
		
	\end{itemize}
	
\end{idea}
	
	\begin{idea}[Логика доказательства]
		
		Самые важные переходы:
		
		\begin{itemize}
			
			\item \textbf{(2) $\Rightarrow$ (1) (От плотности к сходимости ряда):}
			
			Мы знаем, что к вектору $h$ можно приблизиться \textit{какой-то} линейной комбинацией (по определению плотности). Но нам нужно доказать, что именно \textit{ряд Фурье} сходится к $h$.
			
			Здесь работает \textbf{Предложение \ref{minproperty}}: оно гарантирует, что ряд Фурье приближает \textit{лучше всех}. Если какая-то комбинация дает ошибку $\varepsilon$, то ряд Фурье даст ошибку $\le \varepsilon$. Значит, он сходится.
			
			\item \textbf{(1) $\Leftrightarrow$ (3) (Базис $\Leftrightarrow$ Парсеваль):}
			
			Это следует из тождества:
			
			\[ \|\text{Ошибка приближения}\|^2 = \|h\|^2 - \sum |c_k|^2 \]
			
			Если ошибка стремится к 0 (базис), то сумма квадратов стремится к норме (Парсеваль), и наоборот.
			
			
			
			\item \textbf{(2) $\Leftrightarrow$ (4) (Плотность $\Leftrightarrow$ Перпендикуляры):}
			
			Это геометрический факт: пространство $H$ раскладывается на замыкание оболочки $\overline{[e]}$ и её ортогональное дополнение. Чтобы $\overline{[e]}$ было всем пространством, дополнение должно быть нулевым.
			
		\end{itemize}
		
	\end{idea}

\begin{proof}
	\begin{enumerate}
		\item Так как $H$ сепарабельно, в нем существует счетное всюду плотное множество $F = \{f_n\}$ (т.е. $\overline{[F]} = H$).
		\item Применим к $\{f_n\}$ процесс ортогонализации Грама-Шмидта. Получим ортонормированную систему $e = \{e_n\}$.
		\item Процесс сохраняет линейные оболочки: $[e_1, \dots, e_n] = [f_1, \dots, f_n]$. Значит $[e] = [F]$.
		\item Следовательно $\overline{[e]} = \overline{[F]} = H$. По теореме \ref{thm4.4}, система $e$ является базисом.
	\end{enumerate}
\end{proof}

\begin{remark}
	Это означает, что все сепарабельные гильбертовы пространства одной размерности изоморфны (как векторные пространства со скалярным произведением). Например, $L_2[0, 1]$ изоморфно пространству последовательностей $l_2$.
\end{remark}